\section{Ubicación de la clase: no se trata simplemente de «hacer menos profundo el Transformer»}

El término «casi superficial» (near-shallow) del título puede hacer creer que la clase solo trata de reducir el número de capas del Transformer. En realidad, el ponente propone una ruta más concreta y más sujeta a condiciones. Primero, reescribe una parte del trabajo de «proyectar primero y comparar después» de la self-attention estándar como «transformar directamente las comparaciones entre vectores»; segundo, usa la generalized co-occurrence (coocurrencia generalizada) para construir un warm start no aleatorio para la softmax layer; tercero, cuando el embedding está fijo, almacena en caché las comparaciones entre vectores de la capa inferior, de modo que una parte del cálculo ya no se repite en cada entrenamiento o inferencia; cuarto, observa si esta estructura permite un tiny controller entrenado solo con datos locales de la aplicación. Aquí, near-shallow significa que \textbf{el cálculo reutilizable de la capa inferior hace menos profunda la ruta efectiva en línea}, no que a la red le quede una sola capa.

\subsection{Primero se corrigen las fuentes: el título y el contenido de las notas anteriores no eran confiables}

La versión anterior fechaba la clase en 2026, usaba una URL de una Stanford lecture inexistente y añadía mucho contenido como drift gate, attention observability board, fairness dashboard, weekly newsletter y RLHF governance roadmap. Estos mecanismos no aparecen en la página oficial del curso ni en la grabación, ni tampoco en los artículos relacionados del ponente. Esta clase toma ahora como única base la sesión del 2 de mayo de 2024, la grabación oficial de Stanford Online \texttt{zL9B3eXq0gY}, los subtítulos oficiales \texttt{en-US} hechos por personas y los 27 slide states didácticos recuperados de la grabación.

\lecturefigure{slide-01-motivation.jpg}{Pregunta central: ¿se puede aprender directamente la transformación de «resultados de comparación a feature weights»?}{Grabación oficial de Stanford Online 00:04:39}

La primera slide didáctica compara dos rutas. La attention estándar cambia el espacio de representación mediante $W_q,W_k$ y luego produce el score con el dot product; la forma alternativa del ponente conserva la comparación entre los vectores originales $XX^\top$ y usa otra matriz $W$ para convertir esas comparaciones en attention logits. Toda la inicialización explícita, el almacenamiento en caché y los experimentos de edge que siguen giran en torno a esta reescritura.

\begin{importantbox}{Las tres preguntas centrales de esta clase}
\begin{enumerate}
  \item \textbf{Initialization:} ¿se puede evitar arrancar con parámetros aleatorios y obtener un valor inicial útil directamente a partir de estadísticas de entrada y salida?
  \item \textbf{Reusable computation:} ¿qué representaciones o comparaciones entre vectores pueden congelarse y almacenarse en caché para reducir el cálculo repetido?
  \item \textbf{Application-only learning:} ¿puede un modelo pequeño aprender solo de los datos de interacción de una tarea local, sin depender de un preentrenamiento general?
\end{enumerate}
\end{importantbox}

\teachervoice{00:01:19--00:02:49, Williams explica primero su formación en matemáticas, física y quantitative linguistics. Parte deliberadamente de la estructura estadística y de fórmulas explícitas, y no de las scaling recipes habituales.}

\subsection{Niveles de evidencia: el conocimiento estándar, las derivaciones de los artículos y los experimentos de clase no deben mezclarse}

En esta clase aparecen a la vez fórmulas estándar del Transformer, derivaciones de preprints del equipo, resultados experimentales en las slides de clase, una demostración de viabilidad en Le Potato y planteamientos a futuro en la sesión de Q\&A. Para no presentar un prototipo de investigación como una receta consolidada, conviene mantener la siguiente separación por niveles al leer.

\begin{knowledgebox}{Registro de evidencia}
\begin{tabularx}{\textwidth}{L{0.17\textwidth} L{0.29\textwidth} X}
\toprule
Nivel & Ejemplos & Lectura correcta \\
\midrule
Antecedentes estándar & softmax, dot-product attention, perplejidad, padding & Puede tomarse como base técnica común \\
Derivaciones de artículos & solución explícita de generalized co-occurrence, SAFFU hidden targets & Deben respetarse los supuestos y las condiciones de aproximación del artículo \\
Experimentos de clase & curvas warm/cold, MNIST, BabyLM, switch task & Solo respaldan lo observado en la configuración correspondiente \\
Juicios del ponente & el packing no es el context model correcto, nombre de PLM & Son posturas de investigación, no consenso \\
Trabajo futuro & convolution warm start, multimodal, RLHF, implementación pública & Aún no está terminado; no debe presentarse como capacidad verificada \\
\bottomrule
\end{tabularx}
\end{knowledgebox}

\begin{warningbox}{SAFFU no es una arquitectura estándar disponible}
Self-Attentive Feed-Forward Unit (SAFFU, unidad de prealimentación autoatenta) y Precision Language Model (PLM, modelo de lenguaje de precisión) son términos de investigación que usa el equipo del ponente. No son módulos estandarizados en PyTorch/Hugging Face, y tampoco puede concluirse, solo con base en esta clase, que «entrenar un LLM ya no requiere backpropagation».
\end{warningbox}

\subsection{Resumen de la sección}

La clase real gira en torno a la inicialización no aleatoria, el almacenamiento en caché y los modelos locales pequeños, no en torno a los procesos de gobernanza de la versión anterior. Estas notas ubicarán cada resultado dentro de sus supuestos concretos: softmax, características no negativas, embedding fijo, datos limitados y un prototipo de edge en etapa temprana.
